% ── Marco Teórico ────────────────────────────────────────────────────────────
\section{Marco Teórico}

\subsection{Arquitectura Cliente-Servidor y API REST}
El sistema sigue una arquitectura cliente-servidor desacoplada, en la que el
cliente (frontend) y el servidor (backend) se comunican exclusivamente a
través de una interfaz de programación de aplicaciones (API) que expone
recursos mediante el protocolo HTTP siguiendo los principios REST
(Representational State Transfer): uso de verbos HTTP (\texttt{GET},
\texttt{POST}, \texttt{PUT}, \texttt{PATCH}, \texttt{DELETE}) sobre recursos
identificados por rutas, intercambio de datos en formato JSON y ausencia de
estado en el servidor entre peticiones, delegando la sesión del usuario al
token enviado en cada solicitud.

\subsection{Node.js y Express}
Node.js es un entorno de ejecución de JavaScript del lado del servidor basado
en el motor V8, orientado a operaciones de entrada/salida no bloqueantes.
Express es un framework minimalista para Node.js que provee un sistema de
enrutamiento y de middlewares para construir APIs HTTP \parencite{expressdocs}.
En el presente proyecto, Express se utiliza para definir los módulos de rutas,
controladores y middlewares que conforman la API REST del backend.

\subsection{PostgreSQL}
PostgreSQL es un sistema de gestión de bases de datos relacionales de código
abierto que implementa el estándar SQL con soporte para transacciones ACID,
llaves foráneas, restricciones de verificación (\texttt{CHECK}) e índices
\parencite{postgresqldocs}. El modelo de datos del sistema se apoya en estas
características para garantizar la integridad referencial entre solicitudes,
usuarios, trámites y documentos.

\subsection{React y Vite}
React es una biblioteca de JavaScript para la construcción de interfaces de
usuario basada en componentes reutilizables y un modelo declarativo de
renderizado \parencite{reactdocs}. Vite es la herramienta de compilación y
servidor de desarrollo utilizada para el proyecto frontend, que ofrece
recarga en caliente y un empaquetado optimizado para producción.

\subsection{Autenticación basada en JSON Web Tokens (JWT)}
JWT es un estándar abierto (RFC 7519) que define un formato compacto y
autocontenido para transmitir información entre dos partes como un objeto
JSON firmado digitalmente \parencite{rfc7519}. En el sistema, el token JWT se
emite tras un inicio de sesión exitoso y contiene el identificador del
usuario y su rol, permitiendo validar cada solicitud subsecuente sin
necesidad de mantener sesiones en el servidor.

\subsection{Hash de Contraseñas}
Para el almacenamiento seguro de credenciales, las contraseñas no se guardan
en texto plano, sino como el resultado de una función de hash con salteo
(\emph{salt}) mediante el algoritmo \texttt{bcrypt}, diseñado específicamente
para dificultar ataques de fuerza bruta mediante un factor de costo
configurable.

\subsection{Metodología de Desarrollo}
El desarrollo del proyecto se organizó de forma iterativa e incremental,
construyendo primero el modelo de datos, luego la capa de API y finalmente la
integración con la interfaz de usuario. La evidencia verificable en el
repositorio corresponde a una estructura modular de trabajo por carpetas
(\texttt{backend/} y \texttt{frontend/}), con definición formal de rutas y
servicios y con un flujo de validación local mediante ejecución del servidor y
pruebas automatizadas. No se encuentra evidencia comprobable en el proyecto de
un tablero Kanban o de una herramienta de gestión de sprints distinta a la
organización académica y de entrega del curso, por lo que el enfoque se
describe como una metodología de desarrollo ágil informal, apoyada en la
implementación continua y en pruebas de integración HTTP.

\subsection{Multer y gestión de archivos}
La carga de documentos y comprobantes de pago se implementa en el backend
mediante la librería \texttt{multer}, un middleware de Express para la
recepción de archivos multipart/form-data \\parencite{multerdocs}. En el
repositorio, esta capa se configura con almacenamiento en disco y validación de
tipo de archivo para aceptar únicamente PDF, JPG y PNG, con límite de tamaño de
5 MB por carga.
